% GENERATED FILE. Edit canonical lessons, then run npm run generate:books.
% canonical-source-hash: fnv1a64:b1d7734d5e1174aa
% canonical-lessons: BN-C359-ojuhat, BN-C359-bhan, BN-C359-koiphiyot, BN-C359-pash-katano, BN-C359-grohonjogyo

\chapter{Excuses}
\label{ch:bn-excuses}
\hlchaptermodality{359}

\begin{chapteropening}
\textbf{By the end of this chapter:} \emph{I can give an excuse and tell when one does not hold up.}

Complete the last lesson of chapter 359: I can give an excuse and tell when one does not hold up.
\end{chapteropening}

\section[\texorpdfstring{\bn{অজুহাত} (ôjuhāt)}{অজুহাত (ôjuhāt)}]{\bn{অজুহাত} (ôjuhāt) — an excuse; a pretext}

\label{lesson:BN-C359-ojuhat}



\begin{quote}
Before the new one: say the Bengali for ordinary, then the Bengali for to mix up.
\end{quote}

\subsection*{You'll want to know: \bn{অজুহাত}}
\textbf{\bn{অজুহাত}} — \emph{ôjuhāt} — \textquotedblleft{}an excuse; a pretext\textquotedblright{}.

\begin{quote}
\bn{দেরির} \bn{অজুহাত} \bn{দিও} \bn{না।}
\end{quote}

\begin{quote}
\emph{derir ôjuhāt dio nā.}
\end{quote}

\begin{quote}
Don't make excuses for being late.
\end{quote}

\begin{cousinweb}
\textbf{\bn{অজুহাত}} (\emph{ôjuhāt}) is usually traced to Arabic \emph{wujūhāt}, a plural meaning reasons or aspects, through Persian: a reason offered to get out of something.
\end{cousinweb}

\begin{grammarlens}[title={What you've built}]
One more word for giving reasons and opinions.
\end{grammarlens}

\subsection*{Your turn}
\begin{itemize}
\raggedright
  \item \emph{Say it:} \emph{ôjuhāt}
  \item \emph{Say it:} \emph{ôjuhāt}, once more
  \item \emph{Say it:} \emph{derir ôjuhāt dio nā.} — Don't make excuses for being late.
  \item \emph{Recall:} say the Bengali for ordinary, then the Bengali for to mix up, then say \emph{ôjuhāt} again
\end{itemize}

\begin{tcolorbox}[breakable,colback=teal!4,colframe=teal!35!black,title={Before you move on}]
What does \bn{অজুহাত} mean? (\textquotedblleft{}An excuse; a pretext\textquotedblright{}.) Say it once more.
\end{tcolorbox}
\section[\texorpdfstring{\bn{ভান} (bhān)}{ভান (bhān)}]{\bn{ভান} (bhān) — a pretence}

\label{lesson:BN-C359-bhan}



\begin{quote}
Before the new one: say the Bengali for to mix up, then the Bengali for an excuse.
\end{quote}

\subsection*{You'll want to know: \bn{ভান}}
\textbf{\bn{ভান}} — \emph{bhān} — \textquotedblleft{}a pretence\textquotedblright{}.

\begin{quote}
\bn{অরুণ} \bn{না} \bn{শোনার} \bn{ভান} \bn{করল।}
\end{quote}

\begin{quote}
\emph{ôrun nā shonār bhān kôrlô.}
\end{quote}

\begin{quote}
Arun pretended not to hear.
\end{quote}

\begin{cousinweb}
\textbf{\bn{ভান}} (\emph{bhān}) is Sanskrit \emph{bhāna}, appearance (Monier-Williams; CDIAL 9451 lists Bengali bhān). Bengali uses it for a false appearance: \bn{ভান} \bn{করা}, to pretend.
\end{cousinweb}

\begin{grammarlens}[title={What you've built}]
One more word for giving reasons and opinions.
\end{grammarlens}

\subsection*{Your turn}
\begin{itemize}
\raggedright
  \item \emph{Say it:} \emph{bhān}
  \item \emph{Say it:} \emph{bhān}, once more
  \item \emph{Say it:} \emph{ôrun nā shonār bhān kôrlô.} — Arun pretended not to hear.
  \item \emph{Recall:} say the Bengali for to mix up, then the Bengali for an excuse, then say \emph{bhān} again
\end{itemize}

\begin{tcolorbox}[breakable,colback=teal!4,colframe=teal!35!black,title={Before you move on}]
What does \bn{ভান} mean? (\textquotedblleft{}A pretence\textquotedblright{}.) Say it once more.
\end{tcolorbox}
\section[\texorpdfstring{\bn{কৈফিয়ত} (koiphiyôt)}{কৈফিয়ত (koiphiyôt)}]{\bn{কৈফিয়ত} (koiphiyôt) — a demanded explanation; an account of oneself}

\label{lesson:BN-C359-koiphiyot}



\begin{quote}
Before the new one: say the Bengali for an excuse, then the Bengali for a pretence.
\end{quote}

\subsection*{You'll want to know: \bn{কৈফিয়ত}}
\textbf{\bn{কৈফিয়ত}} — \emph{koiphiyôt} — \textquotedblleft{}a demanded explanation; an account of oneself\textquotedblright{}.

\begin{quote}
\bn{ম্যানেজার} \bn{দেরির} \bn{জন্য} \bn{কৈফিয়ত} \bn{চাইলেন।}
\end{quote}

\begin{quote}
\emph{mænejār derir jônyô koiphiyôt chāilen.}
\end{quote}

\begin{quote}
The manager asked for an explanation for the delay.
\end{quote}

\begin{cousinweb}
\textbf{\bn{কৈফিয়ত}} (\emph{koiphiyôt}) is Persian \emph{kaifiyat}, how things stand, an account, usually traced to Arabic \emph{kayfiyya}. In Bengali it is the account you are asked to give of yourself.
\end{cousinweb}

\begin{grammarlens}[title={What you've built}]
One more word for giving reasons and opinions.
\end{grammarlens}

\subsection*{Your turn}
\begin{itemize}
\raggedright
  \item \emph{Say it:} \emph{koiphiyôt}
  \item \emph{Say it:} \emph{koiphiyôt}, once more
  \item \emph{Say it:} \emph{mænejār derir jônyô koiphiyôt chāilen.} — The manager asked for an explanation for the delay.
  \item \emph{Recall:} say the Bengali for an excuse, then the Bengali for a pretence, then say \emph{koiphiyôt} again
\end{itemize}

\begin{tcolorbox}[breakable,colback=teal!4,colframe=teal!35!black,title={Before you move on}]
What does \bn{কৈফিয়ত} mean? (\textquotedblleft{}A demanded explanation; an account of oneself\textquotedblright{}.) Say it once more.
\end{tcolorbox}
\section[\texorpdfstring{\bn{পাশ} \bn{কাটানো} (pāsh kāṭāno)}{পাশ কাটানো (pāsh kāṭāno)}]{\bn{পাশ} \bn{কাটানো} (pāsh kāṭāno) — to sidestep; to dodge}

\label{lesson:BN-C359-pash-katano}



\begin{quote}
Before the new one: say the Bengali for a pretence, then the Bengali for a demanded explanation.
\end{quote}

\subsection*{You'll want to know: \bn{পাশ} \bn{কাটানো}}
\textbf{\bn{পাশ} \bn{কাটানো}} — \emph{pāsh kāṭāno} — \textquotedblleft{}to sidestep; to dodge\textquotedblright{}.

\begin{quote}
\bn{মন্ত্রী} \bn{প্রশ্নটা} \bn{পাশ} \bn{কাটালেন।}
\end{quote}

\begin{quote}
\emph{môntrī prôshnoṭā pāsh kāṭālen.}
\end{quote}

\begin{quote}
The minister sidestepped the question.
\end{quote}

Past (he/she …): \bn{তিনি} \bn{বিষয়টা} \bn{পাশ} \bn{কাটালেন}, \emph{tini bishôyṭā pāsh kāṭālen}.

\begin{cousinweb}
\textbf{\bn{পাশ} \bn{কাটানো}} (\emph{pāsh kāṭāno}) is \bn{পাশ}, the side (CDIAL 8118 lists Bengali pās 'near'), with \bn{কাটানো}, to cut: to cut past at the side, to go round a question instead of through it.
\end{cousinweb}

\begin{grammarlens}[title={What you've built}]
One more word for giving reasons and opinions.
\end{grammarlens}

\subsection*{Your turn}
\begin{itemize}
\raggedright
  \item \emph{Say it:} \emph{pāsh kāṭāno}
  \item \emph{Say it:} \emph{pāsh kāṭāno}, once more
  \item \emph{Say it:} \emph{môntrī prôshnoṭā pāsh kāṭālen.} — The minister sidestepped the question.
  \item \emph{Recall:} say the Bengali for a pretence, then the Bengali for a demanded explanation, then say \emph{pāsh kāṭāno} again
\end{itemize}

\begin{tcolorbox}[breakable,colback=teal!4,colframe=teal!35!black,title={Before you move on}]
What does \bn{পাশ} \bn{কাটানো} mean? (\textquotedblleft{}To sidestep; to dodge\textquotedblright{}.) Say it once more.
\end{tcolorbox}
\section[\texorpdfstring{\bn{গ্রহণযোগ্য} (grôhônjogyô)}{গ্রহণযোগ্য (grôhônjogyô)}]{\bn{গ্রহণযোগ্য} (grôhônjogyô) — acceptable}

\label{lesson:BN-C359-grohonjogyo}



\begin{quote}
Before the new one: say the Bengali for a demanded explanation, then the Bengali for to sidestep.
\end{quote}

\subsection*{You'll want to know: \bn{গ্রহণযোগ্য}}
\textbf{\bn{গ্রহণযোগ্য}} — \emph{grôhônjogyô} — \textquotedblleft{}acceptable\textquotedblright{}.

\begin{quote}
\bn{তোমার} \bn{অজুহাতটা} \bn{গ্রহণযোগ্য} \bn{নয়।}
\end{quote}

\begin{quote}
\emph{tomār ôjuhātṭā grôhônjogyô nôy.}
\end{quote}

\begin{quote}
Your excuse is not acceptable.
\end{quote}

\begin{cousinweb}
\textbf{\bn{গ্রহণযোগ্য}} (\emph{grôhônjogyô}) is Sanskrit \emph{grahaṇa}, taking or acceptance, with \emph{yogya}, fit for (Monier-Williams): fit to be accepted. An excuse, a proposal or a reason can be \bn{গ্রহণযোগ্য}; the compound itself is not in the dictionary.
\end{cousinweb}

\begin{grammarlens}[title={What you've built}]
One more word for giving reasons and opinions.
\end{grammarlens}

\subsection*{Your turn}
\begin{itemize}
\raggedright
  \item \emph{Say it:} \emph{grôhônjogyô}
  \item \emph{Say it:} \emph{grôhônjogyô}, once more
  \item \emph{Say it:} \emph{tomār ôjuhātṭā grôhônjogyô nôy.} — Your excuse is not acceptable.
  \item \emph{Recall:} say the Bengali for a demanded explanation, then the Bengali for to sidestep, then say \emph{grôhônjogyô} again
\end{itemize}

\begin{tcolorbox}[breakable,colback=teal!4,colframe=teal!35!black,title={Before you move on}]
What does \bn{গ্রহণযোগ্য} mean? (\textquotedblleft{}Acceptable\textquotedblright{}.) Say it once more.
\end{tcolorbox}
